\documentclass{ximera}

\input{../../../preamble.tex}


\definecolor{fvdnavy}{HTML}{071D43}
\definecolor{fvdcyan}{HTML}{20BFC6}
\definecolor{fvdcyanlight}{HTML}{EFFCFB}
\definecolor{fvdmagenta}{HTML}{F10D69}
\definecolor{fvdmagentalight}{HTML}{FFF1F7}
\definecolor{fvdtext}{HTML}{163044}





%=================================================
% Pedagogische omgevingen
% PDF: tcolorbox
% HTML: learning-box
%=================================================

\newenvironment{voorkennis}
{%
    \ifdefined\HCode
        \HCode{%
<section class="learning-box learning-box--prior">
<div class="learning-box__header">
<span class="learning-box__icon" aria-hidden="true"></span>
<span class="learning-box__title">Voorkennis</span>
</div>
<div class="learning-box__content">
        }%
        \begin{itemize}
    \else
       \begin{tcolorbox}[
    enhanced,
    breakable,
    colback=fvdcyanlight,
    colframe=fvdcyan,
    coltext=fvdtext,
    boxrule=0.5pt,
    leftrule=4pt,
    arc=4mm,
    outer arc=4mm,
    left=7mm,
    right=6mm,
    top=5mm,
    bottom=5mm,
    before skip=6mm,
    after skip=6mm,
    title=Voorkennis,
    coltitle=fvdcyan,
    fonttitle=\bfseries\Large,
    attach title to upper={\par\smallskip},
    boxed title style={
        empty,
        left=0mm,
        right=0mm,
        top=0mm,
        bottom=0mm
    },
    overlay={
        \draw[
            fvdcyan,
            line width=0.5pt
        ]
        ([xshift=42mm,yshift=-13mm]frame.north west)
        --
        ([xshift=-7mm,yshift=-13mm]frame.north east);
    }
]
        \begin{itemize}
    \fi
}
{%
    \end{itemize}
    \ifdefined\HCode
        \HCode{%
</div>
</section>
        }%
    \else
        \end{tcolorbox}
    \fi
}


\newenvironment{leerdoelen}
{%
    \ifdefined\HCode
        \HCode{%
<section class="learning-box learning-box--goals">
<div class="learning-box__header">
<span class="learning-box__icon" aria-hidden="true"></span>
<span class="learning-box__title">Leerdoelen</span>
</div>
<div class="learning-box__content">
        }%
        \begin{itemize}
    \else
\begin{tcolorbox}[
    enhanced,
    breakable,
    colback=fvdmagentalight,
    colframe=fvdmagenta,
    coltext=fvdtext,
    boxrule=0.5pt,
    leftrule=4pt,
    arc=4mm,
    outer arc=4mm,
    left=7mm,
    right=6mm,
    top=5mm,
    bottom=5mm,
    before skip=6mm,
    after skip=6mm,
    title=Leerdoelen,
    coltitle=fvdmagenta,
    fonttitle=\bfseries\Large,
    attach title to upper={\par\smallskip},
    boxed title style={
        empty,
        left=0mm,
        right=0mm,
        top=0mm,
        bottom=0mm
    },
    overlay={
        \draw[
            fvdmagenta,
            line width=0.5pt
        ]
        ([xshift=42mm,yshift=-13mm]frame.north west)
        --
        ([xshift=-7mm,yshift=-13mm]frame.north east);
    }
]
        \begin{itemize}
    \fi
}
{%
    \end{itemize}
    \ifdefined\HCode
        \HCode{%
</div>
</section>
        }%
    \else
        \end{tcolorbox}
    \fi
}


\addPrintStyle{../../..}

\title{Oefeningen — De inverse functie}
\author{Ingmar Herreman}

\begin{document}

% FVD_CHAPTER_DATA_START
% Automatisch gegenereerd uit index.tex.
% Niet handmatig aanpassen.
\fvdchapterdata{4}{Van relatie naar functie}{4}
% FVD_CHAPTER_DATA_END







\maketitle
\FVDbeginExercises



\begin{oefening}[Vooruit en terug]{\oefB\ESS}

Een bijectie \(f\) voldoet aan
\[
f(2)=7,\qquad f(5)=1,\qquad f(-1)=4.
\]

\begin{question}
Bepaal \(f^{-1}(7)\), \(f^{-1}(1)\) en \(f^{-1}(4)\).

\begin{oplossing}
\[
f^{-1}(7)=2,\qquad
f^{-1}(1)=5,\qquad
f^{-1}(4)=-1.
\]
\end{oplossing}
\end{question}

\begin{question}
Welk punt op de grafiek van \(f^{-1}\) hoort bij het punt \((2,7)\) op \(f\)?

\begin{oplossing}
\[
(7,2).
\]
\end{oplossing}
\end{question}

\end{oefening}

\begin{oefening}[Domein en bereik wisselen]{\oefB\ESS}

Gegeven is een bijectie \(f\) met
\[
\operatorname{dom}f=[-3,5]
\]
en
\[
\operatorname{ber}f=[2,10].
\]

\begin{question}
Bepaal \(\operatorname{dom}f^{-1}\).

\begin{oplossing}
\[
\operatorname{dom}f^{-1}=[2,10].
\]
\end{oplossing}
\end{question}

\begin{question}
Bepaal \(\operatorname{ber}f^{-1}\).

\begin{oplossing}
\[
\operatorname{ber}f^{-1}=[-3,5].
\]
\end{oplossing}
\end{question}

\end{oefening}

\begin{oefening}[Inverse grafiek]{\oefB\ESS}

Een bijectieve functie gaat door de punten
\[
(-2,1),\quad(0,2),\quad(3,5),\quad(4,7).
\]

\begin{question}
Geef vier punten van de inverse functie.

\begin{oplossing}
\[
(1,-2),\quad(2,0),\quad(5,3),\quad(7,4).
\]
\end{oplossing}
\end{question}

\begin{question}
Leg uit welke meetkundige transformatie de punten omzet.

\begin{oplossing}
De coördinaten worden verwisseld. Dat is een spiegeling ten opzichte van
de rechte \(y=x\).
\end{oplossing}
\end{question}

\end{oefening}

\begin{oefening}[Inverse van een lineaire functie]{\oefB\ESS}

Bepaal telkens de inverse functie.

\begin{question}
\[
f(x)=2x+5.
\]

\begin{oplossing}
\[
y=2x+5
\]
geeft na verwisselen
\[
x=2y+5.
\]
Dus
\[
\boxed{f^{-1}(x)=\frac{x-5}{2}.}
\]
\end{oplossing}
\end{question}

\begin{question}
\[
g(x)=\frac{x-4}{3}.
\]

\begin{oplossing}
\[
y=\frac{x-4}{3}
\Longrightarrow
x=3y+4.
\]
Na verwisselen volgt
\[
\boxed{g^{-1}(x)=3x+4.}
\]
\end{oplossing}
\end{question}

\begin{question}
Controleer één van beide antwoorden door de functies na elkaar uit te voeren.

\begin{oplossing}
Bijvoorbeeld:
\[
f^{-1}(f(x))
=
\frac{(2x+5)-5}{2}
=x.
\]
\end{oplossing}
\end{question}

\end{oefening}

\begin{oefening}[Kwadraatfunctie met domeinbeperking]{\oefB\ESS}

Gegeven
\[
f:[0,+\infty[\to[4,+\infty[:x\mapsto x^2+4.
\]

\begin{question}
Bepaal \(f^{-1}\).

\begin{oplossing}
\[
y=x^2+4
\]
en dus na verwisselen
\[
x=y^2+4.
\]
Daaruit:
\[
y=\sqrt{x-4}.
\]
Dus
\[
\boxed{f^{-1}(x)=\sqrt{x-4}.}
\]
\end{oplossing}
\end{question}

\begin{question}
Geef domein en bereik van \(f^{-1}\).

\begin{oplossing}
\[
\operatorname{dom}f^{-1}=[4,+\infty[,
\qquad
\operatorname{ber}f^{-1}=[0,+\infty[.
\]
\end{oplossing}
\end{question}

\end{oefening}




\begin{oefening}[Grafiek en inverse samen analyseren]{\oefU\ESS}

Beschouw
\[
f(x)=\frac12x+2.
\]

\begin{question}
Bepaal algebraïsch \(f^{-1}\).

\begin{oplossing}
\[
y=\frac12x+2
\Longrightarrow
x=\frac12y+2
\Longrightarrow
y=2x-4.
\]
Dus
\[
f^{-1}(x)=2x-4.
\]
\end{oplossing}
\end{question}

\begin{question}
Teken \(f\), \(f^{-1}\) en \(y=x\) in hetzelfde assenstelsel.

\begin{oplossing}
De rechten \(y=\frac12x+2\) en \(y=2x-4\) zijn elkaars spiegelbeeld
ten opzichte van \(y=x\).
\end{oplossing}
\end{question}

\begin{question}
Bepaal het snijpunt van \(f\) met \(f^{-1}\).

\begin{oplossing}
Een punt dat bij spiegeling in \(y=x\) op zichzelf blijft, ligt op \(y=x\).
Dus:
\[
x=\frac12x+2
\Longrightarrow x=4.
\]
Het snijpunt is
\[
(4,4).
\]
\end{oplossing}
\end{question}

\end{oefening}

\begin{oefening}[GeoGebra: ontdek de spiegeling]{\oefU\ESS}

Voer in GeoGebra
\[
f(x)=3x-2
\]
in.

\begin{enumerate}
    \item Bepaal zelf \(f^{-1}\) en voer die functie ook in.
    \item Teken \(y=x\).
    \item Kies een beweegbaar punt \(A\) op \(f\).
    \item Spiegel \(A\) in \(y=x\).
    \item Onderzoek of het spiegelpunt steeds op \(f^{-1}\) ligt.
    \item Formuleer een verklaring met coördinaten.
\end{enumerate}

\begin{oplossing}
\[
f^{-1}(x)=\frac{x+2}{3}.
\]
Als
\[
A=(a,f(a))=(a,b),
\]
dan is het spiegelpunt
\[
A'=(b,a).
\]
Omdat \(f^{-1}(b)=a\), ligt \(A'\) op de grafiek van \(f^{-1}\).
\end{oplossing}

\end{oefening}

\begin{oefening}[Welke inverse hoort erbij?]{\oefU\ESS}

Beschouw
\[
f:[0,+\infty[\to[2,+\infty[:x\mapsto (x+1)^2+1.
\]

\begin{question}
Bepaal de inverse functie.

\begin{oplossing}
\[
y=(x+1)^2+1.
\]
Verwissel:
\[
x=(y+1)^2+1.
\]
Dus
\[
x-1=(y+1)^2.
\]
Omdat \(y\geq0\) in het bereik van de inverse:
\[
y+1=\sqrt{x-1}.
\]
Daarom
\[
\boxed{f^{-1}(x)=\sqrt{x-1}-1.}
\]
\end{oplossing}
\end{question}

\begin{question}
Waarom nemen we niet de negatieve wortel?

\begin{oplossing}
Omdat het bereik van \(f^{-1}\) gelijk is aan het domein van \(f\),
namelijk \([0,+\infty[\). De negatieve tak zou niet met die domeinkeuze
overeenkomen.
\end{oplossing}
\end{question}

\end{oefening}

\begin{oefening}[Geen inverse functie zonder beperking]{\oefU\ESS}

Beschouw
\[
f:\mathbb R\to[0,+\infty[:x\mapsto(x-2)^2.
\]

\begin{question}
Waarom bestaat er op dit domein geen inverse functie?

\begin{oplossing}
De functie is niet injectief. Bijvoorbeeld
\[
f(1)=f(3)=1.
\]
De functiewaarde \(1\) heeft dus twee originelen.
\end{oplossing}
\end{question}

\begin{question}
Geef twee mogelijke domeinbeperkingen die een inverse functie mogelijk maken.

\begin{oplossing}
\[
[2,+\infty[
\qquad\text{of}\qquad
]-\infty,2].
\]
\end{oplossing}
\end{question}

\begin{question}
Bepaal voor beide keuzes de inverse.

\begin{oplossing}
Voor domein \([2,+\infty[\):
\[
f^{-1}(x)=2+\sqrt{x}.
\]

Voor domein \(]-\infty,2]\):
\[
f^{-1}(x)=2-\sqrt{x}.
\]
\end{oplossing}
\end{question}

\end{oefening}

\begin{oefening}[Inverse in een sensorcontext]{\oefU\ESS}

Een sensor heeft het kalibratiemodel
\[
U(T)=0{,}40+0{,}025T,
\qquad 0\leq T\leq100.
\]

\begin{question}
Bepaal het bereik van \(U\).

\begin{oplossing}
\[
U(0)=0{,}40,\qquad U(100)=2{,}90.
\]
Dus
\[
\operatorname{ber}U=[0{,}40,2{,}90].
\]
\end{oplossing}
\end{question}

\begin{question}
Bepaal de inverse omzetting \(T(U)\).

\begin{oplossing}
\[
U=0{,}40+0{,}025T
\Longrightarrow
\boxed{T(U)=\frac{U-0{,}40}{0{,}025}.}
\]
\end{oplossing}
\end{question}

\begin{question}
Een meettoestel registreert \(1{,}65\) V. Welke temperatuur rapporteert het model?

\begin{oplossing}
\[
T(1{,}65)=\frac{1{,}65-0{,}40}{0{,}025}=50.
\]
Dus \(50^\circ\mathrm C\).
\end{oplossing}
\end{question}

\begin{question}
Leg in woorden uit wat de oorspronkelijke functie en de inverse functie
elk doen.

\begin{oplossing}
De oorspronkelijke functie zet temperatuur om in sensorspanning.
De inverse functie zet een gemeten sensorspanning terug om in de
bijbehorende temperatuur.
\end{oplossing}
\end{question}

\end{oefening}




\begin{oefening}[Snijpunten van een functie en haar inverse]{\oefG}

Een leerling beweert:

\begin{center}
``Een functie en haar inverse snijden elkaar altijd op de rechte \(y=x\).''
\end{center}

Onderzoek de uitspraak kritisch.

\begin{question}
Leg uit waarom punten die door de spiegeling in \(y=x\) op zichzelf blijven,
op \(y=x\) liggen.

\begin{oplossing}
Een punt \((a,b)\) blijft bij het verwisselen van coördinaten hetzelfde
precies wanneer \(a=b\). Zulke punten liggen op \(y=x\).
\end{oplossing}
\end{question}

\begin{question}
Betekent dit automatisch dat elke functie en haar inverse een snijpunt hebben?

\begin{oplossing}
Nee. Een functie kan bijectief zijn zonder een vast punt \(f(x)=x\)
in haar domein te hebben. De grafieken zijn wel elkaars spiegelbeeld,
maar hoeven elkaar niet noodzakelijk te snijden.
\end{oplossing}
\end{question}

\begin{question}
Welke vergelijking zoek je als je gemeenschappelijke punten op \(y=x\)
wil vinden?

\begin{oplossing}
\[
f(x)=x.
\]
\end{oplossing}
\end{question}

\end{oefening}

\begin{oefening}[Ontwerp een omkeerbare meetfunctie]{\oefG}

Een meetsysteem moet een grootheid \(q\in[0,10]\) coderen als een spanning
\(U\in[1,5]\). Het systeem moet zo ontworpen zijn dat \(q\) uit \(U\)
eenduidig kan worden teruggevonden.

\begin{enumerate}
    \item Ontwerp een eenvoudige lineaire bijectie \(U(q)\).
    \item Bepaal de inverse \(q(U)\).
    \item Controleer beide samenstellingen.
    \item Teken beide grafieken samen met \(y=x\).
    \item Leg uit waarom bijectiviteit voor zo'n meetsysteem inhoudelijk noodzakelijk is.
\end{enumerate}

\begin{oplossing}
Een natuurlijke keuze is de rechte door \((0,1)\) en \((10,5)\):
\[
U(q)=1+0{,}4q.
\]
Daaruit:
\[
q(U)=\frac{U-1}{0{,}4}.
\]
De samenstellingen geven telkens de oorspronkelijke invoer terug.
Bijectiviteit garandeert dat elke toegelaten sensoruitvoer bij precies één
waarde van de gemeten grootheid hoort.
\end{oplossing}

\end{oefening}

\begin{oefening}[Voorbereiding op cyclometrische functies]{\oefG}

Beschouw de grafiek van
\[
f(x)=\sin x.
\]

Gebruik GeoGebra.

\begin{enumerate}
    \item Onderzoek met horizontale lijnen waarom sinus op \(\mathbb R\)
          niet injectief is.
    \item Zoek een interval waarop iedere waarde tussen \(-1\) en \(1\)
          precies één keer wordt aangenomen.
    \item Leg uit waarom een domeinbeperking noodzakelijk is als we later
          een inverse van sinus willen definiëren.
\end{enumerate}

\begin{oplossing}
Sinus is periodiek en neemt dezelfde waarden oneindig vaak aan.
Een standaard geschikt interval is
\[
\left[-\frac{\pi}{2},\frac{\pi}{2}\right].
\]
Op dit interval neemt sinus iedere waarde uit \([-1,1]\) precies eenmaal aan.
Met dit beperkte domein ontstaat een bijectie naar \([-1,1]\), zodat een
inverse functie kan worden gedefinieerd. De volledige cyclometrische functie
wordt later systematisch behandeld.
\end{oplossing}

\end{oefening}





\end{document}
